\documentclass[12pt, a4paper]{article}
\usepackage[english,spanish,es-tabla ]{babel}
\usepackage[top=0.5cm, left=0.5cm, right=0.5cm, bottom=0.5cm]{geometry}
\usepackage{mystyle_2}
\usepackage{mycolours}
\usepackage{mytcolorbox}
\usepackage{mycommands}
\usepackage{microtype}
\usepackage[all]{chemmacros}

\usepackage{pdflscape}
\usepackage{rotating}
\usepackage{tabularray}
\UseTblrLibrary{siunitx}
\newcommand{\acu}{{(aq)}}

\begin{document}

\DefTblrTemplate{contfoot-text}{normal}{Continúa en la siguiente página}
\SetTblrTemplate{contfoot-text}{normal}
\DefTblrTemplate{conthead-text}{normal}{(Continuación)}
\SetTblrTemplate{conthead-text}{normal}

\begin{landscape}

    \begin{longtblr}[caption=Compuestos binarios]{hlines,vlines, colspec={Q[c,m] Q[c,m, 2cm] Q[c,m,7cm] Q[c,m,7cm]}, rowsep=5pt, rowhead=1}
        \SetRow{gray!50} nº  & Fórmula  & Tradicional  & Hidrógeno  \\
        1 & \ch{HBrO}  & \color{red} ácido hipobromoso &\color{red} hidrogeno(oxidobromato)\\
        2 & \ch{HBrO2}  & \color{red}ácido bromoso  &\color{red}  hidrogeno(dioxidobromato)\\
        3 & \ch{HBrO3}  & \color{red} ácido brómico &\color{red} hidrogeno(trioxidobromato) \\
        4 & \ch{HBrO4}  & \color{red} ácido perbrómico &\color{red} hidrogeno(tetraoxidobromato)\\
        5 & \ch{HIO}  & \color{red} ácido hipoyodoso  &\color{red} hidrogeno(oxidoyodato)\\
        6 & \ch{HIO2}  & \color{red} ácido yodoso &\color{red} hidrogeno(dioxidoyodato) \\
        7 & \ch{HIO3}  & \color{red} ácido yódico &\color{red} hidrogeno(trioxidoyodato)\\
        8 & \ch{HIO4}  &\color{red} ácido peryódico  &\color{red} hidrogeno(tetraoxidoyodato)\\
        9 & \ch{HNO2}  &\color{red} ácido nitroso  &\color{red} hidrogeno(dioxidonitrato)\\
        10 & \ch{H2SeO2}  & \color{red} ácido hiposelenoso &\color{red} dihidrogeno(dioxidoselenato)\\
        11 & \ch{H2SO3}  & \color{red} ácido sulfuroso &\color{red} dihidrogeno(trioxidosulfato) \\
        12 & \ch{H2TeO2}  &\color{red}  ácido hipoteluroso  &\color{red} dihidrogeno(dioxidotelurato) \\
        13 & \ch{H2SO4}  & \color{red} ácido sulfúrico & \color{red} dihidrogeno(tetraoxidosulfato)\\
        14 & \ch{H2TeO3}  & \color{red} ácido teluroso & \color{red}dihidrogeno(trixodiotelurato)\\
        15 & \ch{HClO}  & \color{red} ácido hipocloroso & \color{red} hidrogeno(oxidoclorato)\\
        16 & \ch{HClO2}  &  \color{red} ácido cloroso & \color{red}hidrogeno(dioxidoclorato)\\
        17 & \ch{HClO3}  & \color{red} ácido clórico &\color{red} hidrogeno(trioxidoclorato) \\
        18 & \ch{HClO4}  & \color{red} ácido hiperclórico &\color{red} hidrogeno(tetraoxidoclorato)\\
        19 & \ch{HNO3}  & \color{red} ácido nítrico & \color{red} hidrogeno(trioxidonitrato)\\
        20 & \ch{HNO2}  & \color{red} ácido nitroso & \color{red} hidrogeno(dioxidonitrato)\\
        21 & \ch{H2SO2}  & \color{red} ácido hiposulfuroso  &\color{red} dihidrogeno(dioxidosulfato)\\
        22 & \ch{H2SO3}  & \color{red} ácido sulfuroso &\color{red} dihidrogeno(trioxidosulfato)\\
        23 & \ch{H2SO4}  &\color{red} ácido sulfúrico  & \color{red} dihidrogeno(tetraoxidosulfato)\\
        24 & \ch{HMnO4}  & \color{red}  ácido permangánico &\color{red} hidrogeno(tetraoxidomanganato)\\
        25 & \ch{H3PO3}  & \color{red} ácido fosforoso & \color{red} trihidrogeno(trioxidofosfato)\\
        26 & \ch{H2CrO3}  & \color{red} ácido crómico &\color{red} hidrogeno(trioxidocromato)\\
        27 & \ch{H3PO4}  & \color{red} ácido fosfórico &\color{red} trihidrogeno(tetraoxidofosfato)\\
        28 & \ch{H3AsO4}  & \color{red} ácido arsénico & \color{red} trihidrogeno(tetraoxidoarsenato)\\
        29 & \ch{HSbO2}  & \color{red} ácido metaantimonioso & \color{red} hidrogeno(dioxidoantimoniato)\\
        30 & \ch{H3SbO3}  &\color{red} ácido antimonioso ácido &\color{red} trihidrogeno(trioxidoantimoniato)\\
        31 & \ch{HSbO3}  &\color{red}  ácido metaantimónico  &\color{red} hidrogeno(trioxidoantimoniato)\\
        32 & \ch{H3SbO4}  & \color{red}  ácido antimónico&\color{red} trihidrogeno(tetraoxidoantimoniato)\\
        33 & \ch{HNO}  &\color{red} ácido hiponitroso  &\color{red} hidrogeno(oxidonitrato)\\
        34 & \ch{H2CO3}  &\color{red}  ácido carbónico &\color{red} dihidrogeno(trioxidocarbonato)\\
        35 & \ch{H2CO2}  & \color{red} ácido carbonoso & \color{red} dihidrogeno(dioxidocarbonato)\\
        36 & \ch{H2SeO4}  &\color{red} ácido selénico &\color{red} dihidrogeno(tetraoxidoselenato)\\
        37 & \ch{HMnO4}  &\color{red}  ácido permangánico &\color{red} hidrogeno(tetraoxidomaganato)\\
        38 & \ch{H2MnO4}  &\rojo{ácido mangánico}   &\rojo{dihidrogeno(tetraoxidomanganato)} \\
        39 & \ch{H2MnO3}  &\rojo{ácido manganoso}   & \rojo{dihidrogeno(trioxidomanganato)}\\
        40 & \color{red} \ch{HClO}  & ácido hipocloroso  &\color{red} hidrogeno(oxidoclorato)\\
        41 & \color{red} \ch{HIO}  & ácido hipyodoso  &\color{red} hidrogeno(oxidoyodato)\\
        42 & \color{red} \ch{HNO3}  & ácido nítrico  &\color{red} hidrogeno(trioxidonitrato)\\
        43 & \color{red} \ch{H2SeO4}  &  ácido selénico &\color{red} dihidrogeno(tetraoxidoselenato)\\
        44 & \color{red} \ch{H2SO4}  &  ácido sulfúrico &\color{red} dihidrogeno(tetraoxidosulfato)\\
        45 & \color{red} \ch{H2TeO4}  &  ácido telúrico &\color{red} dihidrogeno(tetraoxidotelurato)\\
        46 & \color{red} \ch{HClO2}  & ácido cloroso  &\color{red} hidrogeno(dioxidoclorato)\\
        47 & \color{red} \ch{HClO4}  & ácido perclórico  &\color{red} hidrogeno(tetraoxidoclorato)\\
        48 & \color{red} \ch{HClO}  & ácido hipocloroso  &\color{red} hidrogeno(oxidoclorato)\\
        49 & \color{red} \ch{HClO3}  &  ácido clórico &\color{red} hidrogeno(trioxidoclorato)\\
        50 & \color{red} \ch{HBrO}  &  ácido hipobromoso  &\color{red} hidrogeno(oxidobromato)\\
        51 & \color{red} \ch{HBrO2}  &  ácido bromoso  &\color{red} hidrogeno(dioxidobromato)\\
        52 & \color{red} \ch{HBrO3}  &  ácido brómico &\color{red} hidrogeno(trioxidobromato)\\
        53 & \color{red} \ch{HBrO4}  & ácido perbrómico   &\color{red} hidrogeno(tetraoxidobromato) \\
        54 & \color{red} \ch{H2MnO4}  &  ácido mangánico  &\color{red} dihidrogeno(tetraoxidomanganato)\\
        55 & \color{red} \ch{HMnO4}  &  ácido permangánico  &\color{red} hidrogeno(tetraoxidomanganato)\\
        56 & \color{red} \ch{H2CrO4}  &  ácido crómico  &\color{red} dihidrogeno(tetraoxidocromato)\\
        57 & \color{red} \ch{H2SO2}  &  ácido hiposulfuroso  &\color{red} dihidrogeno(dioxidosulfato)\\
        58 & \color{red} \ch{H2SO4}  &  ácido sulfúrico   &\color{red} dihidrogeno(tetraoxidosulfato)\\
        59 & \color{red} \ch{H2TeO4}  &  ácido telúrico  &\color{red} dihidrogeno(tetraoxidotelurato)\\
        60 & \color{red} \ch{HNO2}  & ácido nitroso   &\color{red} hidrogeno(dioxidonitrato)\\
        61 & \color{red} \ch{H3PO4}  &  ácido fosfórico  &\color{red} trihidrogeno(tetraoxidofosfato)\\
        62 & \color{red} \ch{HNO3}  &   ácido nítrico &\color{red} hidrogeno(trioxidonitrato)\\
        63 & \color{red} \ch{H2SO3}  &  ácido sulfuroso  &\color{red} dihidrogeno(trioxidosulfato)\\
        64 & \color{red} \ch{H2TeO3}  &  ácido teluroso  &\color{red} dihidrogeno(trioxidotelurato)\\
        65 & \color{red} \ch{H2SeO4}  &  ácido selénico  &\color{red} dihidrogeno(tetraoxidoselenato)\\
        66 & \color{red} \ch{H2CO3}  &  ácido carbónico  &\color{red} dihidrogeno(trioxidocarbonato)\\
        67 & \color{red} \ch{H2CO2}  &  ácido carbonoso  &\color{red} dihidrogeno(dioxidocarbonato)\\
        68 & \color{red} \ch{H4SiO4}  &  ácido silícico  &\color{red} tetrahidrogeno(tetraoxidosilicato)\\
        69 & \color{red} \ch{H3AsO3}  &   ácido arsenicoso &\color{red} trihidrogeno(trioxidoarsenato)\\
        70 & \color{red} \ch{HBrO2}  &  \color{red} ácido bromoso & Hidrogeno(dioxidobromato)\\
        71 & \color{red} \ch{HBrO}  & \color{red} ácido hipobromoso  & Hidrogeno(oxidobromato) \\
        72 & \color{red} \ch{HBrO3}  &  \color{red} ácido brómico &  Hidrogeno(trioxidobromato) \\
        73 & \color{red} \ch{HBrO4}  &  \color{red} ácido perbrómico & Hidrogeno(tetraoxidobromato) \\
        74 & \color{red} \ch{HIO}  &  \color{red}  ácido hipoyodoso & Hidrogeno(oxidoyodato)\\
        75 & \color{red} \ch{HIO2}  &  \color{red} ácido yodoso & Hidrogeno(dioxidoyoadato)\\
        76 & \color{red} \ch{HIO3}  &  \color{red} ácido yódico  &  Hidrogeno(trioxidoyodato)\\
        77 & \color{red} \ch{HIO4}  &  \color{red} ácido peryódico & Hidrogeno(tetraoxidoyodato)\\
        78 & \color{red} \ch{HNO2}  &  \color{red} ácido nitroso & Hidrogeno(dioxidonitrato)\\
        79 & \color{red} \ch{H2SeO2}  & \color{red} ácido hiposelenoso  &Dihidrogeno(dioxidoselenato) \\
        80 & \color{red} \ch{H2SO3}  &  \color{red}  ácido sulfuroso &Dihidrogeno(trioxidosulfato)\\
        81 & \color{red} \ch{H2SO4}  &  \color{red} ácido sulfúrico &Dihidrogeno(tetraoxidosulfato) \\
        82 & \color{red} \ch{H2TeO3}  &  \color{red} ácido teluroso &Dihidrogeno(trioxidotelurato) \\
        83 & \color{red} \ch{H2TeO4}  &  \color{red} ácido telúrico &Dihidrogeno(tetraoxidotelurato) \\
        84 & \color{red} \ch{HClO}  &  \color{red} ácido hipocloroso &Hidrogeno(oxidoclorato) \\
        85 & \color{red} \ch{HClO2}  &  \color{red} ácido clororo &Hidrogeno(dioxidoclorato) \\
        86 & \color{red} \ch{HClO3}  &  \color{red}  ácido clórico &Hidrogeno(trioxidoclorato) \\
        87 & \color{red} \ch{HBrO4}  &  \color{red} ácido brómico &Hidrogeno(tetraoxidobromato) \\
        88 & \color{red} \ch{HBrO}  &  \color{red} ácido hipobromoso &Hidrogeno(oxidobromato) \\
        89 & \color{red} \ch{HBrO2}  &  \color{red} ácido bromoso  &Hidrogeno(dioxidobromato) \\
        90 & \color{red} \ch{H2SeO4}  &  \color{red} ácido selénico &Dihidrogeno(tetraoxidoselenato) \\
        91 & \color{red} \ch{H2SO4}  &  \color{red} ácido sulfúrico &Dihidrogeno(tetraoxidosulfato)\\
        92 & \color{red} \ch{H2TeO4}  &  \color{red} ácido telúrico &Dihidrogeno(tetraoxidotelurato) \\
    \end{longtblr}

\end{landscape}
\end{document}